\section{Multi-Vertical Use Case \& Validation}
\label{chap:proposal:use-case-validation}

To assess the testbed's suitability for the development and validation of vertical applications in an end-to-end scenario, a Network Application will need to be developed and evaluated in the proposed testbed. This will demonstrate the testbed's capabilities and showcase how its underlying infrastructure and interfaces can be used to deploy and validate applications under the requirements of a specific vertical use case. Section \ref{sota:networkapplications:usecases} highlighted several possible network applications across distinct verticals and scenarios, such as \ac{CAM}, \ac{T&L}, and \ac{PPDR}. Given the testbed's goal of catering to multiple vertical industries, the considered \ac{PoC} targets the \ac{PPDR} and Automotive sectors through a vehicle collision detection and notification solution that bridges the two domains.

This use case is motivated by the need to improve the safety of road users in the moments immediately after an accident. Following a vehicle collision, involved and nearby vehicles and pedestrians remain exposed to the hazards of the crash site, increasing the risk of secondary incidents before emergency services arrive. Mitigating this risk requires detecting incidents quickly and immediately alerting those at and approaching the crash site. In this scenario, a computer vision model would analyze a live video stream from fixed roadside cameras to detect collisions and trigger notifications in real time, requiring close integration with the mobile network to automate detection and alerting.

\hide{
Among these, \acs{PPDR} emerges as the most demanding in terms of service requirements, as failures in communication or service continuity may compromise the outcome of emergency operations, with potentially irreversible consequences. More specifically, scenarios such as \acs{IOPS}~\cite{IOPS5G} and \acs{FIDEGAD}~\cite{FIDEGAD5G} impose strict requirements on communication continuity, latency, and bandwidth, under dynamic and often degraded network conditions. This makes \acs{PPDR} an ideal vertical context for a \acl{PoC}, where a network application operating under such conditions would necessarily have to be deeply integrated with the \acs{QoS}, Monitoring, and Analytics interfaces provided by the testbed to securely adapt to network changes and enforce the service guarantees that mission-critical services demand.

The \acs{PoC} application is designed within the \acs{FIDEGAD} scenario~\cite{FIDEGAD5G}, in which drone-based fire detection and ground assistance capabilities are deployed to support emergency response teams in wildfire situations. Inherently, the envisioned network application consists of two main components: (i) a camera client, intended to be deployed at the drone (i.e. the \acs{UE}), that captures live video and streams it over a \acs{5G} connection; and (ii) a server, intended to be deployed at the \acs{MEC} or the cloud, that receives the stream, processes it through an \acs{AI} fire detection model and presents the results through a web interface accessible to ground operators, raising an alert whenever fire is identified in a frame. Given the mission-critical nature of this scenario, the network application must sustain continuous video delivery and respond dynamically to changes in network conditions, adjusting its \acs{QoS} profile as handovers are performed between cells, which naturally occurs as the drone surveys different geographical areas in search for fire outbreaks.
}

\subsection{Application Interaction Flow}
\label{chap:proposal:use-case-validation:flow}

Figure~\ref{fig:proposal:use-case-validation:flow:overview} illustrates a high-level view of the proposed solution and the interactions between its components. The envisioned system consists of three main elements: (i) a fixed roadside camera, acting as an \acs{UE}, that captures live video and streams it over the \acs{5G} network; (ii) an inference server, deployed at the edge or the cloud, that processes the received streams through an \acs{AI} vehicle collision detection model; and (iii) the actual network application that, through the northbound interfaces of the mobile network, manages the network resources supporting the video stream and alerts nearby road users once a collision is detected. In addition, operator dashboards should provide \acs{PPDR} agents with a real-time overview of the monitored road segment, displaying the live camera feeds alongside the detected collision events so they can inspect the scene and corroborate the outputs produced by the \acs{AI} model. Each event must be georeferenced and represented on an interactive map, together with the real-time positions of nearby road users. This information is also expected to be provided by the network.


\begin{figure}[htbp]
    \centering
    \includegraphics[width=1\linewidth]{figs/placeholder.png}
    \caption{Conceptual overview of the proposed solution}
    \label{fig:proposal:use-case-validation:flow:overview}
\end{figure}



\subsection{Application Requirements}
\label{chap:proposal:use-case-validation:requirements}

With the context and purpose of the network application established, the next step is to derive the requirements for the \acs{PoC} application. These will be organized into two categories: functional and non-functional requirements. Functional requirements describe the system and its functionality, i.e. the tasks it is capable of performing. Non-functional requirements specify the conditions under which the system is required to operate, as well as the properties it must exhibit. This category comprises performance requirements, which define the extent to which the system's functions must be executed through quantitative and individually verifiable metrics, and quality requirements, which capture system criteria such as security and portability.

\subsubsection{Functional Requirements}
\label{chap:proposal:use-case-validation:requirements:functional}

The functional requirements align with the primary functionalities provided by the application. Given the emergency nature of the use case, stakeholders must be able to reserve the network resources needed to support the application. In particular, the uplink connection from the roadside camera to the inference server must satisfy defined throughput and latency requirements to ensure reliable, real-time detection. These requirements should be specified in advance, as part of the application's \acs{QoS} profile, and verified throughout the application's lifecycle. This requires continuous visibility into the network conditions experienced by the application, as the degradation of the communication layer may warrant, for instance, an adjustment of \acs{QoS} to preserve detection performance.

To achieve continuous monitoring and adaptive \acs{QoS}, the application requires exposure interfaces that allow it to interact programmatically with the underlying mobile network and request the necessary network capabilities and \acsp{KPI}. The same applies to other services that are crucial to its operation. The application must identify and track vehicle and pedestrian \acs{UE} within a defined perimeter around the monitored area, to determine which users to alert in case of an emergency. Once a collision is detected, these users must then be notified either individually, e.g. by \acs{MSISDN}, or collectively, by targeting a specific class of \acsp{UE}, such as those associated with pedestrians or vehicles. These capabilities, however, depend on the network infrastructure, and the testbed, exposing the necessary functionality to third-party applications through appropriate interfaces. Accordingly, the functional requirements of the application are defined as follows:

\begin{custenumerate}{FR}
    \item \label{enum:fr:slice} The system shall request, prior to its operation, the reservation of network resources dedicated exclusively to the application, isolated from other traffic sharing the mobile network;
    \item \label{enum:fr:stream} The system shall establish a live video stream from the roadside camera to the inference server over the \acs{5G} network;
    \item \label{enum:fr:detection} The system, i.e. the inference server, shall analyze each frame received from the video stream to detect vehicle collisions;
    \item \label{enum:fr:detection-event} The system shall report each detected vehicle collision as an event from the inference server to the network application;
    \item \label{enum:fr:qos-profile} The system shall define a \acs{QoS} profile for the camera's data session, specifying the uplink throughput and latency to be requested from the network;
    \item \label{enum:fr:monitoring-profile} The system shall define a monitoring profile for the camera's data session, specifying the throughput and latency thresholds against which the observed network conditions should be evaluated;
    \item \label{enum:fr:qos-request} The system shall request the network \acsp{KPI} specified in the active \acs{QoS} profile for the camera's data session;
    \item \label{enum:fr:monitoring} The system shall monitor the network conditions experienced by the camera's data session, comparing them against the thresholds defined in the monitoring profile;
    \item \label{enum:fr:qos-change} The system shall change the active \acs{QoS} profile of the camera's data session whenever the monitored network conditions violate the thresholds defined in the monitoring profile;
    \item \label{enum:fr:ue-tracking} The system shall identify and actively track all vehicle and pedestrian \acsp{UE} located within or entering a predefined area around the surveilled area;
    \item \label{enum:fr:notify-pedestrians} The system shall notify all pedestrian \acsp{UE} located within the predefined area identified in \ref{enum:fr:ue-tracking}, at the moment the collision detection event is processed;
    \item \label{enum:fr:notify-vehicles} The system shall notify all vehicle \acsp{UE} located within the predefined area identified in \ref{enum:fr:ue-tracking}, at the moment the collision detection event is processed;
    \item \label{enum:fr:notify-periodic} The system shall periodically reissue alert notifications to the appropriate \acsp{UE} until the incident is deemed resolved.
\end{custenumerate}

\subsubsection{Non-Functional Requirements}
\label{chap:proposal:use-case-validation:requirements:nonfunctional}

The non-functional requirements specify the operational constraints that govern the execution of the functionalities defined above. Among these, the performance requirements are presented first, each linked to the functional requirements specified in Section~\ref{chap:proposal:use-case-validation:requirements:functional} whose execution it constrains:

\begin{custenumerate}{NFR}
    \item \label{enum:nfr:uplink-bitrate} The uplink connection between the camera and the inference server shall sustain a guaranteed minimum bitrate sufficient to deliver the video stream at a rate of at least 30 \acs{FPS} (\ref{enum:fr:stream}, \ref{enum:fr:qos-profile}, \ref{enum:fr:qos-request});
    \item \label{enum:nfr:frame-processing} The system, i.e. the inference server, shall process at least one out of every three consecutive frames of the 30 \acs{FPS} input stream, such that no more than two consecutive frames are skipped at any given time (\ref{enum:fr:stream}, \ref{enum:fr:detection});
    \item \label{enum:nfr:degradation-detection} The system shall detect a degradation of the network conditions below the requirements declared in the \acs{QoS} profile within 1 second of its occurrence (\ref{enum:fr:monitoring-profile}, \ref{enum:fr:monitoring});
    \item \label{enum:nfr:qos-enforcement} The system shall enforce the new \acs{QoS} profile within 1 second of the degradation being detected (\ref{enum:fr:qos-change});
    \item \label{enum:nfr:notification-time} The system shall notify the affected \acsp{UE} within 1 second of the collision detection event (\ref{enum:fr:detection-event}, \ref{enum:fr:notify-pedestrians}, \ref{enum:fr:notify-vehicles}).
\end{custenumerate}

It should be noted that the verifiability of some of the performance requirements listed above is not determined exclusively by the envisioned system, but also by the capabilities provided by the underlying mobile network. As the core functionality of the application is dependent on the streaming of a live video over \acs{5G}, the requirements on uplink bitrate (\ref{enum:nfr:uplink-bitrate}), and consequently the sustained 30~\acs{FPS} stream, can only be satisfied if the network is configured with a \acs{QoS} profile that defines the necessary resource guarantees, and if the network can actually fulfill it. Similarly, the timing requirements for monitoring (\ref{enum:nfr:degradation-detection}), \acs{QoS} adaptation (\ref{enum:nfr:qos-enforcement}) and notifications (\ref{enum:nfr:notification-time}) partially depend on the processing and response times of the northbound interfaces exposed by the testbed, outside of the developer's responsibility.

The remaining requirements address the quality attributes of the system, which are independent of a specific operational scenario. These are not tied to any individual function and are transversal to the entire system:

\begin{custenumerate}{NFR}
    \setcounter{enumi}{5}
    \item \label{enum:nfr:secure-exposure} All northbound \acs{API} interactions shall be discovered, authenticated and authorized through a single secure exposure layer;
    \item \label{enum:nfr:logging} The system shall log all positive detection events and interactions with network \acsp{API}, including \acs{QoS} and location changes, to support debugging and validation procedures;
    \item \label{enum:nfr:packaging} All application components shall be packaged as \acs{NFV} or containerized applications, deployable through the testbed's orchestration infrastructure.
\end{custenumerate}

\subsection{Interface Requirements}

The functional, performance, and non-functional requirements defined in the previous section characterize the system's behavior and constraints. From these emerge a complementary category of requirements that is particularly relevant to the context of this dissertation, centered on a Testbed for Testing Network Exposure Interfaces: the interface requirements. Interface requirements specify how the system is required to interact with external systems (external interfaces) or with other elements within the system (internal interfaces)~\cite{ISORequirementsEngineering}.

More specifically, defining the interface requirements entails specifying the actual northbound \acsp{API} that will be supporting network applications, as well as any other related components pertaining to the exposure and accessibility of said \acsp{API}. With this goal, Table~\ref{tab:poc:api:mapping} maps each of the application requirements defined at Section \ref{chap:proposal:use-case-validation:requirements} to the network interface or component responsible for fulfilling it, showcasing how the \acs{PoC} integrates with the components described in the preceding sections.

\plan{ Detail
\begin{enumerate}
    \item usually AsSessionWithQoS uses references. Should the C1 constraint be there?
\end{enumerate}
}

\begin{table}[ht]
\centering
\small
\renewcommand{\arraystretch}{1.3}
\rowcolors{2}{gray!15}{white}
\begin{tabular}{p{3.9cm} p{1.4cm} p{8.1cm}}
\toprule
\textbf{Interface} & \textbf{Req.} & \textbf{Role in the \acs{PoC}} \\
\midrule
\parbox[t]{3.9cm}{\acs{NEF} \textit{AsSessionWithQoS}} & FR5, C1, PR1, PR5 & Enforcing \acs{QoS} profile, guaranteeting minimum uplink bitrate for drone/\acs{UE} flow \\
\parbox[t]{3.9cm}{\acs{NEF} \textit{MonitoringEvent}} & FR4, PR4 & Subscribing \acs{UE} location change events (type LOCATION\_REPORTING), notifying the application to trigger a \acs{QoS} update upon drone area transitions \\
\parbox[t]{3.9cm}{\acs{NEF} \textit{AnalyticsExposure}} & PR1 & Monitoring of live throughput, volume and \acs{RTT} for continuous performance validation \\
\mbox{CAMARA \acs{QoS} Profiles} \acs{API} & C1 & Defining the set of available \acs{QoS} profiles and associated parameters, allowing CAMARA \acs{QoD} \acs{API} to select and enforce the requested profile \\
\mbox{CAMARA \acs{QoD} \acs{API}} & FR5, PR1, PR5 & High-level \acs{API} for enforcing the \acs{QoS} profile, guaranteeting minimum uplink bitrate for drone/\acs{UE} flow and abstracting the underlying \acs{NEF} \textit{AsSessionWithQoS} \acs{API} \\
\mbox{CAMARA Geofencing} \mbox{Subscriptions \acs{API}} & FR4, FR5, PR4 & High-level \acs{API} for subscribing to geofencing events that trigger a \acs{QoS} update as the drone entries a defined geographical area, abstracting \acs{NEF} \textit{MonitoringEvent} \acs{API} (type LOCATION\_REPORTING) \\
\mbox{CAMARA Location} \mbox{Retrieval \acs{API}} & FR4 & High-level \acs{API} for retrieving the drone's current location to determine the active operational area and abstracting the underlying \acs{NEF} \textit{MonitoringEvent} \acs{API} (type LOCATION\_REPORTING)  \\
\mbox{CAMARA Location} \mbox{Verification \acs{API}} & FR4 & High-level \acs{API} for verifying whether the drone is within a specified geographical area and abstracting the underlying \acs{NEF} \textit{MonitoringEvent} \acs{API} (type LOCATION\_REPORTING)  \\
\mbox{CAMARA Application} \mbox{Profiles \acs{API}} & C1 & High-level \acs{API} declaring the application's network requirements as a profile, to be used by the Connectivity Insights Subscriptions \acs{API} \\
\mbox{CAMARA Connectivity} \mbox{Insights Subscriptions \acs{API}} & PR1 & High-level \acs{API} for subscribing to network analytics events for validation procedures, abstracting the underlying \acs{NEF} \textit{AnalyticsExposure} \acs{API} \\
\acs{CAPIF} & NR1, NR3 & Onboarding, discovery, authentication and authorization for all testbed \acsp{API}, acting as a single registration point for the network application \\
\bottomrule
\end{tabular}
\caption{Fulfillment of \acs{PoC} application requirements via network exposure interfaces}
\label{tab:poc:api:mapping}
\end{table}

As shown in Table~\ref{tab:poc:api:mapping}, \acs{CAPIF} establishes a single secure exposure layer, through which a variety of \acs{NEF} and CAMARA northbound \acsp{API} are made accessible to network applications. It also reveals that several application requirements, i.e. FR4, FR5, C1, PR1, PR4 and PR5, are satisfied by both \acs{NEF} and CAMARA interfaces simultaneously. However, it is worth emphasizing that the sets of \acs{NEF} and CAMARA \acsp{API} are not redundant, but alternatives targeting different consumer profiles.

On one hand, the \acs{NEF} interfaces are defined by \acs{3GPP} allowing precise control over network capabilities, such as \acs{QoS} session management, event monitoring and analytics exposure. These interfaces expose network functionalities with a degree of technical specificity that requires a deep understanding of the underlying \acs{5G} architecture, data models and overall knowledge of the broader ecosystem.

On the other hand, CAMARA \acsp{API}, aligned with the objectives defined by the CAMARA project~\cite{CAMARAProject}, provide a higher-level user-friendly abstraction, destined to developers with no telco expertise, allowing them to focus on the behavior of their own network applications rather than the complexities of the underlying network, while still taking advantage of its capabilities. Nevertheless, the CAMARA approach still relies on the same underlying \acs{NEF} interfaces to actually deliver network services, according to the information flows described at Section \ref{chap:proposal:exposure:camara}. In the following section, the end-to-end interaction flow of the proposed application is defined, detailing the sequence of operations performed across the various system components and highlighting the difference between the \acs{NEF} and CAMARA \acsp{API} approaches.

\plan{
    Include an interaction sequence diagram illustrating the end-to-end flow of the PoC (prioritizing CAMARA):
    \begin{enumerate}
        \item Onboarding via CAPIF
        \item The creation of an application profile declaring application requirements
        \item The creation of a connectivity insights subscription to be notified when the application requirements are not being met
        \item The change of QoS profile in response to being unable to meet the application requirements or in response to a location change (CAMARA QoD API)
    \end{enumerate}
    Include a sequence diagram prioritizing the NEF path. Showcase the added procedural complexity of the NEF implementation flow compared to the CAMARA APIs, in the prespective of the application developer/end-user.
}

\plan{
    ORIGINAL
    \begin{enumerate}
        \item Introduce the PPDR domain as the vertical context for the PoC, drawing from the use cases discussed in Section \ref{sota:networkapplications:usecases} (IOPS, FIDEGAD). Highlight the particular relevance of PPDR scenarios for exercising network exposure interfaces: they require reliable, low-latency communication under dynamic and potentially degraded network conditions, making QoS enforcement and continuous network monitoring essential, not optional, capabilities.
        \item State the objective of the \acs{PoC} application: a Public Protection and Disaster Relief (PPDR) network application developed within the FIDEGAD scenario, where drone-based fire detection and ground assistance capabilities are leveraged to support emergency response teams in wildfire situations. Reenforce this as an ongoing mission-critical operation, dynamically adapting QoS profile in response to changes in network conditions or operational context.
        \item Derive the application requirements from the use case and divide them between Functional Requirements and non functional requirements. Here we have some examples of what's expected:
        \begin{itemize}
            \item Guaranteed minimum UL/DL throughput for video data transmission;
            \item Bounded RTT and low jitter for real-time command and control;
            \item Continuous monitoring of network performance against these requirements;
            \item Ability to escalate to a higher-priority QoS profile when degradation is detected or predicted.
        \end{itemize}
        \item The implemented exposure interfaces, provided by the CAMARA and NEF northbound \acsp{API}, and secured by \acs{CAPIF}, will serve as the foundation for the development and validation of future network applications, by a variety of developers.
        \item Revisit the use cases from Section \ref{sota:networkapplications:usecases} and corresponding internal details (\acs{NEF}/CAMARA \acs{API} suite) from Section \ref{sec:testingandvalidation:netapps}:
        \begin{itemize}
            \item Dynamic \acs{QoS} adaptation enforced through CAMARA \acs{QoS} \acs{API} in response to either geographical changes obtained with CAMARA Device Location \acs{API} in digital twin / \acs{CAM} scenarios, or to network quality deterioration (or to video quality deterioration in case of video streaming/broadcast).
            \item \acs{PPDR}: \acs{IOPS}, \acs{FIDEGAD}
        \end{itemize}
        \item Map each requirement to the corresponding network interface, showing how the PoC will interact with the implemented components to fulfill its goals:
        \begin{itemize}
            \item CAPIF: onboarding and authenticated discovery of all APIs;
            \item NEF AsSessionWithQoS API: direct QoS profile enforcement per UE flow;
            \item NEF MonitoringEvent API: subscription to UE state events (e.g. location changes) that may warrant a QoS adaptation;
            \item NEF Analytics Exposure API: retrieval of live throughput, RTT, jitter, and packet loss metrics;
            \item CAMARA Application Profiles API: declaration of the application's network requirements as a machine-readable profile;
            \item CAMARA Connectivity Insights Subscriptions API: event-driven notification when network performance falls below declared thresholds.
        \end{itemize}
        \item Clarify that NEF and CAMARA APIs are not redundant but complementary alternatives: the CAMARA APIs provide a higher-level, developer-friendly abstraction over the same underlying NEF capabilities, while the NEF APIs offer finer-grained control. Both paths are exercised in the PoC to validate the full integration stack.
        \item Include an interaction sequence diagram illustrating the end-to-end flow of the PoC (prioritizing CAMARA):
        \begin{enumerate}
            \item (in progress)
            \item Onboarding via CAPIF
            \item the creation of application profile according declaring application requirements (does not need to match exactly the slice requirements)
            \item the creation of connectivity insights subscription to be notified when the application requirements are not being respected
            \item  the change of \acs{QoS} profile in response to being unable to meet the application requirements or in response to a location change (CAMARA QoD API)
        \end{enumerate}
        \item Explain how the \acs{PoC} exercises all implemented interfaces end-to-end: it onboards via CAPIF, consumes CAMARA APIs (Application Profiles to declare its requirements, Connectivity Insights Subscriptions to monitor performance, QoD to enforce them)
        \item Include another sequence diagram this time prioritizing the NEF. Showcase the added complexity of the NEF implementation flow comparatively to CAMARA APIs.
    \end{enumerate}
}
